\documentclass[aspectratio=169,11pt]{beamer}
\newcommand{\footerlabel}{DEMO101 Algorithms \textbullet{} Lecture 3}
% 共用样式 —— 虚构示例课件（DEMO），内容全部自编，不来自任何真实课程。
% 用法：每份 deck 先 \documentclass，再定义 \footerlabel，最后 % 共用样式 —— 虚构示例课件（DEMO），内容全部自编，不来自任何真实课程。
% 用法：每份 deck 先 \documentclass，再定义 \footerlabel，最后 % 共用样式 —— 虚构示例课件（DEMO），内容全部自编，不来自任何真实课程。
% 用法：每份 deck 先 \documentclass，再定义 \footerlabel，最后 \input{_style}。
\usepackage[T1]{fontenc}
\usepackage{lmodern}
\usepackage{amsmath,amssymb}
\usetheme{default}
\setbeamertemplate{navigation symbols}{}
\setbeamercolor{frametitle}{fg=blue!50!black}
\setbeamerfont{frametitle}{series=\bfseries,size=\Large}
\setbeamercolor{title}{fg=blue!50!black}
\setbeamercolor{footline}{fg=black!60}
\setbeamertemplate{footline}{%
  \leavevmode\hbox{%
    \begin{beamercolorbox}[wd=.85\paperwidth,ht=2.5ex,dp=1ex,left,leftskip=1em]{footline}\footnotesize\footerlabel\end{beamercolorbox}%
    \begin{beamercolorbox}[wd=.15\paperwidth,ht=2.5ex,dp=1ex,right,rightskip=1em]{footline}\footnotesize\insertframenumber\end{beamercolorbox}}%
  \vskip0pt}
\setbeamertemplate{itemize items}[circle]
。
\usepackage[T1]{fontenc}
\usepackage{lmodern}
\usepackage{amsmath,amssymb}
\usetheme{default}
\setbeamertemplate{navigation symbols}{}
\setbeamercolor{frametitle}{fg=blue!50!black}
\setbeamerfont{frametitle}{series=\bfseries,size=\Large}
\setbeamercolor{title}{fg=blue!50!black}
\setbeamercolor{footline}{fg=black!60}
\setbeamertemplate{footline}{%
  \leavevmode\hbox{%
    \begin{beamercolorbox}[wd=.85\paperwidth,ht=2.5ex,dp=1ex,left,leftskip=1em]{footline}\footnotesize\footerlabel\end{beamercolorbox}%
    \begin{beamercolorbox}[wd=.15\paperwidth,ht=2.5ex,dp=1ex,right,rightskip=1em]{footline}\footnotesize\insertframenumber\end{beamercolorbox}}%
  \vskip0pt}
\setbeamertemplate{itemize items}[circle]
。
\usepackage[T1]{fontenc}
\usepackage{lmodern}
\usepackage{amsmath,amssymb}
\usetheme{default}
\setbeamertemplate{navigation symbols}{}
\setbeamercolor{frametitle}{fg=blue!50!black}
\setbeamerfont{frametitle}{series=\bfseries,size=\Large}
\setbeamercolor{title}{fg=blue!50!black}
\setbeamercolor{footline}{fg=black!60}
\setbeamertemplate{footline}{%
  \leavevmode\hbox{%
    \begin{beamercolorbox}[wd=.85\paperwidth,ht=2.5ex,dp=1ex,left,leftskip=1em]{footline}\footnotesize\footerlabel\end{beamercolorbox}%
    \begin{beamercolorbox}[wd=.15\paperwidth,ht=2.5ex,dp=1ex,right,rightskip=1em]{footline}\footnotesize\insertframenumber\end{beamercolorbox}}%
  \vskip0pt}
\setbeamertemplate{itemize items}[circle]

\title{Algorithms}
\subtitle{Lecture 3: Insertion Sort}
\author{Demo Lecturer}
\date{Semester 2, 2026}

\begin{document}

\begin{frame}
  \titlepage
\end{frame}
% 封面占了两个编号 —— 之后每页的页脚号 = 物理页 + 1（示例：offset 不一定是 0）
\setcounter{framenumber}{2}

\begin{frame}{The Sorting Problem}
  \begin{description}
    \item[Input] a sequence of $n$ items $A[0..n-1]$
    \item[Output] the same items rearranged so that $A[0] \le A[1] \le \dots \le A[n-1]$
  \end{description}
  \bigskip
  Two properties we will ask about every sorting algorithm:
  \begin{itemize}
    \item \textbf{in-place}: uses only $O(1)$ extra memory beyond the array
    \item \textbf{stable}: equal keys keep their original relative order
  \end{itemize}
\end{frame}

\begin{frame}{The Idea: Sort a Hand of Cards}
  Pick up cards one at a time. The cards already in your hand are always sorted.
  \medskip
  \begin{itemize}
    \item Take the next card
    \item Slide it left past every card that is larger
    \item Drop it into the gap
  \end{itemize}
  \bigskip
  In an array: $A[0..i-1]$ is the sorted ``hand''; $A[i]$ is the next card.
\end{frame}

\begin{frame}[fragile]{Insertion Sort: Pseudocode}
\begin{verbatim}
function InsertionSort(A, n)
    for i <- 1 to n - 1
        v <- A[i]
        j <- i - 1
        while j >= 0 and A[j] > v
            A[j + 1] <- A[j]
            j <- j - 1
        A[j + 1] <- v
\end{verbatim}
  \small The basic operation is the comparison \verb|A[j] > v|.
\end{frame}

\begin{frame}{Trace on $[5, 2, 4, 6, 1, 3]$}
  \begin{center}
  \begin{tabular}{cl}
    \hline
    After round $i$ & Array (sorted prefix before the bar) \\
    \hline
    start & $5 \mid 2\; 4\; 6\; 1\; 3$ \\
    $i=1$ & $2\; 5 \mid 4\; 6\; 1\; 3$ \\
    $i=2$ & $2\; 4\; 5 \mid 6\; 1\; 3$ \\
    $i=3$ & $2\; 4\; 5\; 6 \mid 1\; 3$ \\
    $i=4$ & $1\; 2\; 4\; 5\; 6 \mid 3$ \\
    $i=5$ & $1\; 2\; 3\; 4\; 5\; 6$ \\
    \hline
  \end{tabular}
  \end{center}
\end{frame}

\begin{frame}{Why It Is Correct}
  \begin{block}{Loop invariant}
    At the start of the round for index $i$, $A[0..i-1]$ contains the original first $i$
    items in sorted order.
  \end{block}
  \begin{itemize}
    \item \textbf{Initially} ($i = 1$): a single item is sorted
    \item \textbf{Each round}: the inner loop shifts larger items right and inserts $v$, so $A[0..i]$ is sorted
    \item \textbf{At the end} ($i = n$): the whole array is sorted
  \end{itemize}
\end{frame}

\begin{frame}{Complexity and Properties}
  \begin{center}
  \begin{tabular}{ll}
    \hline
    Case & Comparisons \\
    \hline
    Best (already sorted) & $n - 1 \in \Theta(n)$ \\
    Worst (reverse sorted) & $\frac{n(n-1)}{2} \in \Theta(n^2)$ \\
    \hline
  \end{tabular}
  \end{center}
  \medskip
  In-place: yes. Stable: yes (we stop shifting at $A[j] \le v$, never at equal keys).
\end{frame}

\begin{frame}{Summary and Next Time}
  \begin{itemize}
    \item Insertion sort grows a sorted prefix; $\Theta(n^2)$ worst case, $\Theta(n)$ best case
    \item Good on small or nearly sorted inputs
    \item Next lecture: divide-and-conquer and mergesort
  \end{itemize}
\end{frame}

\end{document}
